\par\Needspace{6\baselineskip}\medskip
\begin{center}\hypertarget{native-170-article123}{}\pdfbookmark[2]{123.}{article-123}123.\end{center}
In praec. supposuimus, valores approximatos incognitarum $x,y$ alicunde
iam haberi. Quoties quidem totius orbitae dimensiones approximatae in potestate
sunt (ex aliis forte observationibus per calculos anteriores deductae
iamque per novas corrigendae), conditioni illi absque difficultate satisfieri poterit,
quamcunque significationem incognitis tribuamus. Contra in determinatione
prima orbitae penitus adhuc ignotae (quae est problema longe difficillimum)
neutiquam indifferens est, quasnam incognitas adhibeamus; arte potius
talique modo eligendae sunt, ut valores approximatos ex ipsius problematis
natura haurire liceat. Quod exoptatissime succedit, quoties tres observationes
ad orbitae investigationem adhibitae motum heliocentricum corporis coelestis
non nimis magnum complectuntur. Huiusmodi itaque observationes ad determinationem
primam semper adhibendae sunt, quam dein per observationes
magis ab invicem remotas ad lubitum corrigere conveniet. Nullo enim negotio
perspicitur, observationum errores inevitabiles calculum eo magis turbare,
quo propiores observationes adhibeantur. Hinc colligitur, observationes
ad determinationem primam haud temere eligendas, sed cavendum esse, \emph{primo}
ne sint nimis sibi invicem vicinae, \emph{dein} vero etiam ne nimis ab invicem distent:
in primo enim casu calculus elementorum observationibus satisfacientium expeditissime
quidem absolveretur, sed his elementis ipsis parum fidendum foret,
quinimo erroribus tam enormiter depravata evadere possent, ut ne approximationis
quidem vice fungi valerent; in casu altero vero artificiis, quibus ad determinationem
approximatam incognitarum utendum est, destitueremur, neque
\SourceLeaf{171}inde aliam derivaremus, nisi vel crassissimam ubi hypotheses multo plures, vel
omnino ineptam, ubi tentamina fastidiosissima haud evitare liceret. Sed de
hisce methodi limitibus scite iudicare melius per usum frequentem quam per
praecepta ediscitur: exempla infra tradenda ostendent, ex observationibus Iunonis
22 tantum diebus ab invicem dissitis motumque heliocentricum $7^0 35'$
complectentibus elementa multa iam praecisione gaudentia derivari, ac vicissim,
methodum nostram optimo etiamnum successu ad observationes Cereris applicari,
quae 260 diebus ab invicem distant, motumque heliocentricum $62^0 55'$
includunt, quatuorque hypothesibus seu potius approximationibus successivis
adhibitis elementa optime cum observationibus conspirantia producere.
